\section*{Introduction}
\addcontentsline{toc}{section}{Introduction}
Ce document présente ma note d'intention pour le projet "Ingénieur Citoyen". J'ai choisi de m'engager dans une démarche solidaire en participant à \textbf{La Féminine de Pau} le dimanche 19 avril 2026. Je représenterai à cette occasion mon entreprise d'alternance, \textbf{Safran Landing Systems Bidos}, afin d'allier cohésion professionnelle et engagement caritatif.

\newpage

\section*{Chapitre 1 : Nature et Pertinence du Projet}
\addcontentsline{toc}{section}{Chapitre 1 : Nature et Pertinence du Projet}

\subsection*{La Féminine de Pau : Un engagement caritatif concret}
\phantomsection
\addcontentsline{toc}{subsection}{La Féminine de Pau : Un engagement caritatif concret}
La Féminine de Pau est bien plus qu'une épreuve sportive ; c'est un levier de soutien pour les associations locales de lutte contre le cancer. 
\begin{itemize}
    \item \textbf{Soutien ciblé :} L'édition de printemps à laquelle je participe soutient spécifiquement la lutte contre les maladies inflammatoires de l'intestin via l'association \textbf{AFA} (maladie de Crohn). 
    \item \textbf{Impact local :} Contrairement à des causes nationales diffuses, cet événement aide directement des femmes et des enfants de notre région. Depuis 2012, plus de 80 000 euros ont été reversés grâce aux participants.
    \item \textbf{Lien avec l'ingénieur citoyen :} Cet engagement répond au critère de "contribution au bien commun" et de "lutte contre l'exclusion" par le soutien aux personnes malades.
\end{itemize}

\subsection*{Pertinence au regard des objectifs du CESI}
\phantomsection
\addcontentsline{toc}{subsection}{Pertinence au regard des objectifs du CESI}
Ma participation sous les couleurs de \textbf{Safran Landing Systems Bidos} permet de :
\begin{itemize}
    \item \textbf{Promouvoir la RSE :} Incarner les valeurs de solidarité de mon entreprise sur le territoire palois.
    \item \textbf{Rayonnement :} Valoriser l'image de l'apprenti ingénieur CESI comme un acteur impliqué dans la cité.
\end{itemize}

\subsection*{Conditions de réalisation}
\phantomsection
\addcontentsline{toc}{subsection}{Conditions de réalisation}
Conformément à la note pédagogique, cette action est réalisée sur mon \textbf{temps personnel} (dimanche). Les parties prenantes sont l'association organisatrice (CUP), les bénéficiaires (AFA), Safran et moi-même.

\newpage

\section*{Chapitre 2 : Objectifs SMART et Faisabilité}
\addcontentsline{toc}{section}{Chapitre 2 : Objectifs SMART et Faisabilité}

\subsection*{Définition des objectifs SMART}
\phantomsection
\addcontentsline{toc}{subsection}{Définition des objectifs SMART}
\begin{itemize}
    \item \textbf{Spécifique :} Terminer le parcours de 6 km et contribuer à la collecte de fonds pour l'AFA.
    \item \textbf{Mesurable :} Validation par l'obtention du dossard et le reversement d'une partie de mon inscription à l'association.
    \item \textbf{Atteignable :} Préparation physique progressive débutant 6 semaines avant l'échéance.
    \item \textbf{Réaliste :} Engagement bénévole compatible avec mon rythme d'alternance.
    \item \textbf{Temporellement défini :} Réalisation le dimanche 19 avril 2026.
\end{itemize}

\subsection*{Analyse des risques et Soft Skills}
\phantomsection
\addcontentsline{toc}{subsection}{Analyse des risques et Soft Skills}
Le principal risque est l'aléa physique. Ce projet me permettra de travailler mon \textbf{agilité} et mon \textbf{éthique}, en plaçant l'effort sportif au service d'une cause médicale sérieuse et reconnue.

\newpage

\section*{Conclusion}
\addcontentsline{toc}{section}{Conclusion}
En participant à La Féminine de Pau, je dépasse le simple cadre de l'entreprise pour m'inscrire dans une démarche de citoyenneté active. Ce projet illustre ma volonté d'agir pour la santé publique locale tout en développant mes compétences comportementales d'ingénieur.

\newpage

\section*{Sitographie}
\addcontentsline{toc}{section}{Sitographie}
\begin{itemize}
    \item Site officiel de La Féminine de Pau : \url{https://www.lafemininedepau.fr/}
    \item Association AFA (Lutte contre la maladie de Crohn) : \url{https://www.afa.asso.fr/}
    \item Note pédagogique Ingénieur Citoyen - CESI 2025.
\end{itemize}